\begin{abstract}
Most websites still run PHP applications on version 7.x or older 
that no longer receive official security updates, leaving systems 
exposed to an ever-growing accumulation of known vulnerabilities. 
Manual migration across hundreds of source files is impractical, 
and existing transformation tools do not incorporate security 
compliance checks as an integrated part of the migration workflow. 
Organizations operating under ISO/IEC 27001:2022 must also perform 
security audits as a separate post-migration step, adding to the 
operational burden. This research builds a Python-based pipeline 
that integrates automatic conversion of legacy PHP web applications 
to a configurable target PHP version with ISO/IEC 27001:2022 
secure coding compliance checks in a single, fully local, 
open-source workflow. The pipeline combines Rector for AST-based 
code transformation, Semgrep for security scanning, PHPStan for 
post-conversion type validation, and five open-source large 
language models running locally via Ollama for semantic 
vulnerability analysis and remediation suggestions. Case studies 
were conducted on PHP codebases from FT UGM (245 files) and 
CBT DTI UGM (326 files). Rector achieved conversion rates of 
97.6\% and 91.7\% respectively, with 100\% post-conversion syntax 
validity in both cases. A security finding delta of $\Delta F = 0$ 
confirmed that the conversion process introduced no new 
vulnerabilities. Automatically generated JSON reports include 
per-finding mappings to ISO/IEC 27001:2022 Annex A controls and 
serve directly as audit documentation. Among five evaluated LLMs, 
Qwen2.5-Coder 7b is recommended for pipeline integration, 
consistently achieving a Syntax Validity Rate of 100\% across all 
experimental conditions.
\end{abstract}

\begin{IEEEkeywords}
PHP migration, secure coding, ISO/IEC 27001:2022, large language 
model, AST code transformation
\end{IEEEkeywords}